\scanpage{54}
\S~4 \uline{Topos stratifiant et} \uline{topos cristallin}.

4.1. Interprétation d'une stratification relative d'un Module $\underline F$ sur $X/S$. Si on a une telle stratification, alors pour tt $S$-morphisme $f:\deleted{\unsure{X}}Y\longrightarrow X$, et tt voisinage infinitésimal $Y'$ de $Y$, \inserted{(sur $S$)} \deleted{\unsure{qui qui|qui que|qui puisse}} tel que $f$ veut bien se prolonger (localmt du moins) en $f':Y'\unsure{\mathbin{-}|\longrightarrow}X$, un système \uline{transitif} d'isom entre les $f'{}^*(\underline F)$.
\diagram{\begin{tikzcd}[ampersand replacement=\&,row sep=small]
X \arrow[dd,dash] \& \deleted{\unsure{X}}Y_{\deleted{\unsure{0}}} \arrow[l,"f"'] \arrow[d,arrows={{Hooks[harpoon]}-}] \\
\& \deleted{\unsure{X}}Y'_{\deleted{\unsure{0|n}}} \arrow[ul,"f'"'] \\
S \&
\end{tikzcd}}

On peut donc dire que la \deleted{\unsure{connexion}} \inserted{donnée d'une} stratification \deleted{\unsure{est|une}} \inserted{revient à la} donnée \deleted{\unsure{canonique|un isom}} , pour tt voisinage infinitésimal $Y'$ d'un $Y\unsure{\xrightarrow{f}|\mathrel{\overset f{-}}}X$ (sur $S$), \inserted{soumis à la condition déjà énoncée (\unsure{autom\textsuperscript{t}.|autom.} si $X/S$ form\textsuperscript{t} lisse)} d'un \uline{prolongement} \deleted{\unsure{en}} de $F$ en un \uline{Module} sur $Y'$, -- avec des conditions de fonctorialité évidentes. Pour exprimer plus commodément ce point de vue, et une variante, on va introduire un site ad-hoc.
% archive sheet number: 53

